\documentclass[10pt,a4paper]{amsart}
\usepackage[margin=19mm]{geometry}
\usepackage{amsmath,amssymb,mathtools}
\usepackage[scr=rsfs,cal=euler]{mathalfa}
\usepackage{xfrac}
\usepackage[unicode,hidelinks,bookmarks=false]{hyperref}
\newcommand{\ie}{i.e.\ }
\newcommand{\cf}{cf.\ }
\newcommand{\resp}{respectively\ }
\newcommand{\Subsection}{\subsection}
\providecommand{\nobreakdash}{\nobreak\hbox{-}}
\setlength{\emergencystretch}{2em}
\multlinegap0pt
\usepackage{amssymb}
\usepackage{tikz}
\usepackage[shortlabels]{enumitem}
\usepackage{bm}
\newtheorem{lemma}{Lemma}
\newtheorem{proposition}{Proposition}
\newcommand{\C}{\mathbb C}
\title{On divergence related to Riemann--von Mangoldt's explicit formula of the prime-counting function}
\author{Harald Grobner}
\date{Source: arXiv:2609.02713v1 (CC BY 4.0)}
\begin{document}
\maketitle

\noindent\emph{Source and licence.} On divergence related to Riemann--von Mangoldt's explicit formula of the prime-counting function. Version arXiv:2609.02713v1 (CC BY 4.0), distributed by arXiv under the Creative Commons Attribution 4.0 International licence, http://creativecommons.org/licenses/by/4.0/, as recorded in the arXiv licence metadata for this exact version and shown on the abstract page https://arxiv.org/abs/2609.02713v1. Full licence notice: \texttt{licenses/arxiv-2609.02713-CC-BY-4.0.txt}.\par\smallskip

\noindent\textit{Editorial scope.} Counted unit: the article's proposition prop:HA (source0.tex lines 390--398). The article's complete original proof of it is delivered with it (source0.tex lines 400--446, one \texttt{proof} environment of 47 lines). No mathematics is written, repaired or re-derived: the statement and the proof are the article's own lines, in the article's own notation, and no retained line is substituted or re-flowed.  Retained uncounted as the source-local context the proof needs: lines 348--354.  Nothing was omitted from any retained span, so the package carries no omission record.  No retained line contains a citation, so the wrapper carries no bibliography and no bibliography entry.  No retained line contains a figure environment, an image-inclusion command or drawing code, so no image is delivered and the package declares no asset dependency.  Editorial formatting only, in addition: an AMS \texttt{amsart} preamble that restates the article's own declarations and macros used by the retained lines, namely the environments \texttt{lemma}, \texttt{proposition} on separate counters, together with the standard packages those lines need.  The retained environments print in this document as Lemma 1, Proposition 1 in order of appearance, whereas the article prints the same items with the numbering of its own full document.  The complete native source of the article as distributed in the arXiv e-print for this exact version is bundled unchanged as \texttt{sources/209230/source0.tex}.
\begin{lemma}\label{lem:MA}
The function $M_A$ admits the meromorphic continuation to all $z\in\C$, given by
\begin{equation}\label{eq:MA}
M_A(z)= \sum_{j=0}^{\infty} \frac{(-1)^{j+1}A^j}{j!(z+j)\zeta(z+j+1)}.
\end{equation}
If $\rho$ is a non-trivial zero of $\zeta(s)$, then $M_A(z)$ has pole at $z=\rho-1$.
\end{lemma}

\begin{proposition}\label{prop:HA}
The truncated Mellin transform $H_A$ admits a meromorphic continuation to all $z\in\C$ by
\begin{equation}\label{eq:HAtrans}
H_A(z) = -\Gamma(1-z)\sin\left(\frac{\pi z}{2}\right)\,M_A(z-1) +
\sum_{k=0}^{\infty}
\frac{(-1)^kM_A(2k)}{(2k)!(2k+1-z)}.
\end{equation}
In particular, $H_A$ has a pole at every non-trivial zero $\rho$ of $\zeta(s)$.
\end{proposition}

\begin{proof}
We start off with the following observation: For fixed $u>0$, the integral $\int_1^\infty Y^{-z}\cos(uY)\,dY$ has an analytic continuation to an entire function in $z\in\C$ by virtue of
\begin{equation}\label{eq:KA}
\int_1^\infty Y^{-z}\cos(uY)\,dY=\Gamma(1-z)\sin\left(\frac{\pi z}{2}\right)\,u^{z-1}-\sum_{k=0}^{\infty}\frac{(-1)^k u^{2k}}{(2k)!(2k+1-z)}.
\end{equation}
Indeed, for $0<\Re e(z)<1$, the integral  converges locally uniformly by Dirichlet's test. Moreover, the standard Mellin transform gives
\begin{equation}\label{eq:KA1}
\int_0^\infty Y^{-z}\cos(uY)\,dY = \Gamma(1-z)\sin\left(\frac{\pi z}{2}\right)\,u^{z-1}.
\end{equation}
Integrating the power series of cosine on $[0,1]$ term by term, yields
\begin{equation}\label{eq:KA2}
\int_0^1Y^{-z}\cos(uY)\,dY= \sum_{k=0}^{\infty}\frac{(-1)^k u^{2k}}{(2k)!(2k+1-z)},
\end{equation}
proving \eqref{eq:KA} in the strip $0<\Re e(z)<1$. We claim that the right hand side of \eqref{eq:KA} is an entire function in $z\in\C$: In fact, since $\sin\left(\frac{\pi z}{2}\right)=0$ at positive even integers, the first term $\Gamma(1-z)\sin\left(\frac{\pi z}{2}\right)\,u^{z-1}$ given by \eqref{eq:KA1} is a meromorphic function, whose poles are at the positive odd integers. In order to analyze the second term, given by \eqref{eq:KA2}, let
$$
K\subset\mathbb C\setminus\{2m+1, m\geq 0\}
$$
be any compact subset and choose any $R>0$ such that $|w|\le R$ for all $w\in K$. For all sufficiently large $k$,
$$
|2k+1-w| \ge 2k+1-|w| \ge k \qquad (w\in K).
$$
Hence
$$
\sup_{w\in K} \left| \frac{(-1)^k u^{2k}}{(2k)!(2k+1-w)}\right|\le\frac{u^{2k}}{(2k)!\,k}.
$$
Since $\sum_{k=1}^{\infty} \frac{u^{2k}}{(2k)!\,k}<\infty,$ the series in \eqref{eq:KA2} converges absolutely and uniformly on $K$.  It therefore defines a holomorphic function on $\mathbb C\setminus\{2m+1, m\geq 0\}$, and hence a meromorphic function on $\mathbb C$, with poles (possibly) only at the positive odd integers.\\\\ 
To settle the entireness of the right hand side of \eqref{eq:KA} (and hence our intermediate claim, made at the very beginning of this proof), it remains to show that these apparent poles at the odd positive integers of both terms in \eqref{eq:KA1} and in \eqref{eq:KA2} finally cancel each other in the respective Laurent series expansion, i.e., that in the difference in \eqref{eq:KA} these singularites are removed. For this purpose we return to the original integral.  For $\Re e(z)>0$, integration by parts gives
$$\int_1^\infty Y^{-z}\cos(uY)\,dY=-\frac{\sin u}{u}+\frac{z}{u} \int_1^\infty Y^{-z-1}\sin(uY)\,dY.$$
Indeed, the boundary term at infinity vanishes since $|Y^{-z}\sin(uY)|\le Y^{-\Re e(z)}\longrightarrow 0$ as $Y\to\infty$. The integral on the right is absolutely and locally uniformly convergent in $z$ on the half-plane $\Re e(z)>0$:  More precisely, if $K_0$ is a compact subset of this half-plane, then there exists $\delta>0$ such that $\Re e(w)\ge\delta$ for all $w\in K_0$, and
$$|Y^{-w-1}\sin(uY)| \le Y^{-1-\delta}, \qquad w\in K_0,
$$
and $Y^{-1-\delta}$ is integrable on $[1,\infty)$.  Hence, for fixed $u>0$, $z\mapsto \int_1^\infty Y^{-z}\cos(uY)\,dY$ is in fact holomorphic on $\Re e(z)>0$. The meromorphic expression obtained by the right hand side of \eqref{eq:KA} above agrees with this holomorphic function on the nonempty open strip $0<\Re e(z)<1$, whence by the uniqueness of analytic continuation, the two functions must in fact agree on
$$
\{z\in\C\,| \,\Re e(z)>0\,\}\setminus\{2m+1, m\geq 0\}.
$$
But since $z\mapsto \int_1^\infty Y^{-z}\cos(uY)\,dY$ is holomorphic throughout $\Re e(z)>0$, every apparent singularity of the right hand side of \eqref{eq:KA} at
$z\in \{2m+1, m\geq 0\}$ must vanish in the difference of \eqref{eq:KA1} and \eqref{eq:KA2}. As these were the only possible poles of the right hand side of \eqref{eq:KA} it provides an analytic continuation of $z\mapsto \int_1^\infty Y^{-z}\cos(uY)\,dY$ to an entire function in $z\in\C$ as claimed.\\\\
We may now finish the proof of Prop.\ \ref{prop:HA}: For $\Re e(z)>1$, Fubini applies to the integrals defining $H_A(z)$, because $q_A\in L^1((0,1))$ and, as we have just seen, $\int_1^\infty Y^{-z}\cos(uY)\,dY$ is absolutely convergent. We get
$$
H_A(z)=-\int_0^1q_A(u)\int_1^\infty Y^{-z}\cos(uY)\,dY\,du
$$
and may hence insert the continued expression provided by \eqref{eq:KA}. This yields \eqref{eq:HAtrans} for $\Re e(z)>1$ and hence meromorphic continuation of $H_A(z)$ by Lem.\ \ref{lem:MA}, where we observe that since
$$|M_A(2k)| \leq \int_0^1 |q_A(u)|\,du =\|q_A\|_1<\infty,$$
the series 
$$\sum_{k=0}^{\infty}\frac{(-1)^k M_A(2k)}{(2k)!(2k+1-z)}$$
converges absolutely and locally uniformly on $\mathbb C\setminus\{2m+1, m\geq 0\}$. The last assertion implies furthermore that the series $\sum_{k=0}^{\infty}\frac{(-1)^k M_A(2k)}{(2k)!(2k+1-z)}$ is in fact holomorphic on the whole critical strip and hence on a neighbourhood of every non-trivial zero $\rho$ of $\zeta(s)$. Hence, in order to establish also the last claim about the poles of the right hand side of \eqref{eq:HAtrans}, it suffices to show that the remaining summand $-\Gamma(1-z)\sin\left(\frac{\pi z}{2}\right)\,M_A(z-1)$ in \eqref{eq:HAtrans} has a pole at every non-trivial zero $\rho$ of $\zeta(s)$. But this easily follows, since by Lem.\ \ref{lem:MA}, the function $M_A(z-1)$ has a pole at every non-trivial zero $\rho$ of $\zeta(s)$, whereas $\Gamma(1-\rho)\sin\left(\frac{\pi \rho}{2}\right)\neq 0$.
\end{proof}

\end{document}
